\SourceLeaf{249}%
\[
\begin{aligned}
4EEGGk&=E\mathord{.}\frac{\mathrm{d}E}{\mathrm{d}q}\mathord{.}\frac{\mathrm{d}G}{\mathrm{d}q}+E(\frac{\mathrm{d}G}{\mathrm{d}p})^2+G\mathord{.}\frac{\mathrm{d}E}{\mathrm{d}p}\mathord{.}\frac{\mathrm{d}G}{\mathrm{d}p}\\
&\quad+G(\frac{\mathrm{d}E}{\mathrm{d}q})^2-2EG(\frac{\mathrm{d}\mathrm{d}E}{\mathrm{d}q^2}+\frac{\mathrm{d}\mathrm{d}G}{\mathrm{d}p^2})
\end{aligned}
\]
und für die Änderung des Winkels \(\theta\)
\[
\BareRoot{EG}\mathord{.}\mathrm{d}\theta=\tfrac{1}{2}\mathord{.}\frac{\mathrm{d}E}{\mathrm{d}q}\mathord{.}\mathrm{d}p-\tfrac{1}{2}\mathord{.}\frac{\mathrm{d}G}{\mathrm{d}p}\mathord{.}\mathrm{d}q
\]
\par\indent Unter den verschiedenen Fällen, in denen diese Orthogonalitätsbedingung gilt, ist derjenige der wichtigste, in dem alle Linien eines der beiden Systeme, etwa des ersten, kürzeste Linien sind. Bei einem konstanten Wert von \(q\) ist der Winkel \(\theta\) hier also \(=0\). Die eben angegebene Gleichung für die Änderung des Winkels \(\theta\) zeigt daher, dass \(\frac{\mathrm{d}E}{\mathrm{d}q}=0\) gelten muss, dass also der Koeffizient \(E\) von \(q\) unabhängig sein muss. Somit muss \(E\) entweder konstant oder eine Funktion von \(p\) allein sein. Am einfachsten ist es, für \(p\) die Länge jeder Linie des ersten Systems selbst zu wählen, und zwar von dem Punkt aus gemessen, in dem sich alle Linien des ersten Systems treffen, wenn es einen solchen Punkt gibt, oder, falls sie keinen gemeinsamen Schnittpunkt haben, von einer beliebigen Linie des zweiten Systems aus. Mit diesen Festlegungen ist klar, dass \(p\) und \(q\) nun dieselben Größen bezeichnen, die wir in den Artikeln 15 und 16 durch \(r\) und \(\varphi\) ausgedrückt haben, und dass sich \(E=1\) ergibt. Die beiden vorhergehenden Formeln gehen damit in die folgenden über:
\[
\begin{aligned}
4GGk&=(\frac{\mathrm{d}G}{\mathrm{d}p})^2-2G\frac{\mathrm{d}\mathrm{d}G}{\mathrm{d}p^2}\\
\BareRoot{G}\mathord{.}\mathrm{d}\theta&=-\tfrac{1}{2}\mathord{.}\frac{\mathrm{d}G}{\mathrm{d}p}\mathord{.}\mathrm{d}q
\end{aligned}
\]
oder, wenn wir \(\BareRoot{G}=m\) setzen,
\[
k=-\frac{1}{m}\mathord{.}\frac{\mathrm{d}\mathrm{d}m}{\mathrm{d}p^2},\qquad\mathrm{d}\theta=-\frac{\mathrm{d}m}{\mathrm{d}p}\mathord{.}\mathrm{d}q
\]
Im Allgemeinen ist \(m\) eine Funktion von \(p,q\) und \(m\mathrm{d}q\) der Ausdruck für das Element einer beliebigen Linie des zweiten Systems. Im besonderen Fall, dass alle Linien \(p\) vom selben Punkt ausgehen, muss für \(p=0\) offenbar \(m=0\) gelten. Wählen wir in diesem Fall außerdem für \(q\) den Winkel, den das erste Element einer beliebigen Linie des ersten Systems mit dem Element einer beliebig ausgewählten dieser Linien bildet, so ist für einen unendlich kleinen Wert von \(p\) das Element der Linie des zweiten Systems, die man als einen mit dem Radius \(p\) beschriebenen Kreis auffassen kann, \(=p\mathrm{d}q\). Somit gilt für einen unendlich kleinen Wert von \(p,m=p\). Folglich gelten für \(p=0\) zugleich \(m=0\) und \(\frac{\mathrm{d}m}{\mathrm{d}p}=1\).
\Article{20}
Bleiben wir noch bei derselben Annahme. Es bezeichne also \(p\) allgemein die Länge einer kürzesten Linie, die von einem festen Punkt \(A\) zu einem beliebigen Punkt der Fläche gezogen wird,
